\section*{الفصل \ref{ch:b2:current-potential-curves} --- منحنيات شدة التيار--الجهد}

\begin{solution}{exo:b2:current-potential-curves:1}
يتأكسد الزنك: \omterm{def:b2:current-potential-curves:ie-curve}{تيار مصعدي}، موجب. وتُختزل أيونات النحاس عند قطب
النحاس: \omterm{def:b2:current-potential-curves:ie-curve}{تيار مهبطي}، سالب. والتياران متساويان
بالقيمة المطلقة، لأن الشحنة نفسها تجري في الدارة.
\end{solution}

\begin{solution}{exo:b2:current-potential-curves:2}
$|i_{\lim}| = 96\,485 \times 0.50 \times 10^{-4} \times 1.6 \times 10^{-9} \times
2.0/(25 \times 10^{-6}) = \qty{6.2e-4}{A}$، أي \qty{0.62}{mA}.
\end{solution}

\begin{solution}{exo:b2:current-potential-curves:3}
\omterm{def:b2:current-potential-curves:fast-slow}{النظام السريع} يقطع محور الجهود بحدة عند جهد نرنست له؛ أما \omterm{def:b2:current-potential-curves:fast-slow}{النظام
البطيء} فيُظهر مقطعًا مستويًا من التيار المعدوم حوله، ولا تبدأ فروعه
إلا بعد \omterm{def:b2:current-potential-curves:overpotential}{فرط جهد}. على البلاتين: $\ce{Fe^3+}/\ce{Fe^2+}$ (سريع)؛
$\ce{O2}/\ce{H2O}$ (بطيء).
\end{solution}

\begin{solution}{exo:b2:current-potential-curves:4}
عند جهد نصف الموجة، المساوي $E^\circ = \qty{0.77}{V}$ حين تنتشر الصورتان
بالطريقة نفسها.
\end{solution}

\begin{solution}{exo:b2:current-potential-curves:5}
قرب \qty{0.77}{V} موجة مهبطية لأيونات \ce{Fe^3+} (سريعة)، تبلغ استواءً
يتناسب مع تركيزها. وقرب \qty{0.34}{V} موجة ثانية، إذ يترسب
النحاس؛ ويُضاف استواؤها إلى الأول، وهو تقريبًا ضعفه ارتفاعًا لتركيز
متساوٍ ومعاملي انتشار متقاربين، لأن $n = 2$. ثم، قرب
\qty{0}{V} عند pH 0 على البلاتين (وأدنى منه حين يُغطى القطب بالنحاس)،
الجدار الشديد الانحدار لاختزال \ce{H+}.
\end{solution}

\begin{solution}{exo:b2:current-potential-curves:6}
عند \qty{0.50}{V}: لا يُختزل إلا \ce{Fe^3+} إلى \ce{Fe^2+}، بتياره الحدي.
وعند \qty{0.10}{V}: يُختزل \ce{Fe^3+} ويترسب النحاس في الوقت نفسه،
كل منهما بتياره الحدي؛ والمجموع حاصل جمع الاستواءين.
\end{solution}

\begin{solution}{exo:b2:current-potential-curves:7}
pH 0: \qty{0}{V} و\qty{1.23}{V}؛ وpH 14: \qty{-0.83}{V} و\qty{0.40}{V}. وعلى
البلاتين تحافظ النافذة على عرضها وتنزل بمقدار \qty{59}{mV} لكل وحدة pH،
إذ يتبع كل جدار مزدوجته، ويبقى الجدار المصعدي منزاحًا نحو الأعلى \omterm{def:b2:current-potential-curves:overpotential}{بفرط
جهد} الأكسجين.
\end{solution}

\begin{solution}{exo:b2:current-potential-curves:8}
عند \qty{1.0}{V} يتأكسد \ce{Fe^2+} على استوائه (\omterm{def:b2:current-potential-curves:ie-curve}{تيار مصعدي}
يتناسب مع $[\ce{Fe^2+}]$) ويُختزل كل \ce{Ce^4+} على استوائه
(\omterm{def:b2:current-potential-curves:ie-curve}{تيار مهبطي} يتناسب مع $[\ce{Ce^4+}]$). فقبل التكافؤ ينخفض \omterm{def:b2:current-potential-curves:ie-curve}{التيار
المصعدي} خطيًا مع الحجم المضاف؛ وبعده يكبر \omterm{def:b2:current-potential-curves:ie-curve}{التيار المهبطي}
خطيًا. ويقطع المستقيمان الصفر عند التكافؤ.
\end{solution}

\begin{solution}{exo:b2:current-potential-curves:9}
تتضاعف الاستواءات ($|i_{\lim}| \propto 1/\delta$). ولا يتغير جهد نصف الموجة،
لأن الطبقة واحدة للصورتين؛ ولا يتغير كذلك جهد التيار المعدوم
\omterm{def:b2:current-potential-curves:fast-slow}{لنظام سريع}، وهو جهد نرنست.
\end{solution}

\begin{solution}{exo:b2:current-potential-curves:10}
كما في الفصل مع $i_{\lim,c} = 0$: $E = E_{1/2} + (RT/nF)\ln[i/(i_{\lim} -
i)]$. وباللوغاريتمات العشرية، $E = E_{1/2} + (RT\ln 10/nF)\log[i/(i_{\lim} - i)]$،
وهو مستقيم ميله $0.0592/n$ فولط عند \qty{298}{K}، يقطع
المحور $E$ عند $E_{1/2}$.
\end{solution}

\begin{solution}{exo:b2:current-potential-curves:11}
يستقر الزنك وحده عند الجهد الذي يساوي فيه تياره المصعدي (أكسدة الزنك
السريعة) \omterm{def:b2:current-potential-curves:ie-curve}{التيار المهبطي} لاختزال \ce{H+} على الزنك،
وهو صغير لأن هذا الاختزال بطيء على الزنك: فيكون الذوبان
بطيئًا. وعند ملامسة البلاتين، تستطيع الإلكترونات التي يحررها الزنك اختزال \ce{H+}
على البلاتين، حيث الاختزال سريع: فيكون \omterm{def:b2:current-potential-curves:ie-curve}{التيار المهبطي} عند جهد معيّن
أكبر بكثير، ويُبلغ التوازن عند تيار أعلى، ويذوب
الزنك أسرع بينما يتكوّن الهيدروجين على البلاتين.
\end{solution}

\begin{solution}{exo:b2:current-potential-curves:12}
على البلاتين تبدأ أكسدة الماء قرب \qty{1.6}{V} عند pH 0: فقبل
\qty{2.0}{V} كان التيار سيُصرف في صنع ثنائي الأكسجين. أما القطب الذي
تحتاج عليه أكسدة الماء \omterm{def:b2:current-potential-curves:overpotential}{فرط جهد} أكبر فيدفع الجدار إلى ما بعد
\qty{2.0}{V}، بحيث يمكن أكسدة النوع داخل النافذة.
\end{solution}

\begin{solution}{pb:b2:current-potential-curves:1}
\textbf{1.} \ce{O2 + 2H2O + 4e- -> 4OH-}.
\textbf{2.} \ce{Ag + Cl- -> AgCl + e-} (أربع مرات)؛ والتفاعل الإجمالي \ce{O2 + 2H2O + 4Ag +
4Cl- -> 4AgCl + 4OH-}.
\textbf{3.} $E = 1.229 - 0.0592 \times 7 + (0.0592/4)\log 0.21 = \qty{0.80}{V}$.
\textbf{4.} $E^\circ(\ce{AgCl}/\ce{Ag}) = \qty{0.22}{V}$ مع كلوريد نشاطيته 1:
$-0.60 + 0.22 = \qty{-0.38}{V}$ على سلم الهيدروجين.
\textbf{5.} اختزال \ce{O2} \omterm{def:b2:current-potential-curves:fast-slow}{نظام بطيء}: فتحت جهد نرنست له
يلزم \omterm{def:b2:current-potential-curves:overpotential}{فرط جهد} قدره عدة أعشار الفولط قبل أن يكبر التيار،
وأكثر لبلوغ الاستواء.
\textbf{6.} على الاستواء لا يتعلق التيار إلا بإمداد \ce{O2}،
لا بالانجرافات الصغيرة للجهد: فهو يقيس التركيز.
\textbf{7.} اختزال الماء إلى ثنائي الهيدروجين (الجدار المهبطي)، الذي كان
سيضيف تيارًا لا صلة له بالأكسجين.
\textbf{8.} خطي عبر الغشاء، من $c$ في جهة العيّنة إلى الصفر عند
المهبط.
\textbf{9.} $j = D_m c/L$ لكل وحدة مساحة.
\textbf{10.} $i = -4FAD_m c/L$ (مهبطي)، $|i| = 4FAD_mc/L$.
\textbf{11.} الغشاء هو المرحلة الأبطأ: فالانتشار عبره أبطأ بكثير
منه في الماء خارجه، فيؤدي الغشاء دور \omterm{def:b2:current-potential-curves:limiting-current}{طبقة
الانتشار}.
\textbf{12.} يقع تدرج التركيز داخل الغشاء، الذي لا يتعلق سمكه
بالتحريك (شريطة تجدد الماء خارجه)؛
والترسب على الغشاء يزيد سمكه ويخفض التيار.
\textbf{13.} $0.2095 \times (101.325 - 3.17) = \qty{20.56}{kPa}$.
\textbf{14.} $c = 1.3 \times 10^{-5} \times 20\,560 = \qty{0.27}{mol/m^3}$ ($0.267$).
\textbf{15.} $0.267 \times 32.00 = \qty{8.6}{mg/L}$ (\unit{g/m^3}).
\textbf{16.} $|i| = 4 \times 96\,485 \times 1.0 \times 10^{-5} \times 1.0 \times 10^{-10}
\times 0.267/(25 \times 10^{-6}) = \qty{4.1}{\micro A}$.
\textbf{17.} $4.05/8.55 = \qty{0.474}{\micro A}$ لكل \unit{mg/L}.
\textbf{18.} لا يُعرف سمك الغشاء ومعامل الانتشار فيه بدقة،
ويتغيران مع القِدم ودرجة الحرارة؛ والقياس في محلول معلوم
يتضمنها كلها.
\textbf{19.} صفرًا (وعمليًا تيار متبقٍّ صغير، يُطرح).
\textbf{20.} $2.80/0.474 = \qty{5.9}{mg/L}$.
\textbf{21.} $2.80/4.05 = 69~\%$ من الإشباع.
\textbf{22.} أكسجين تستهلكه عملية تحلل المادة العضوية (مياه الصرف، النباتات الميتة)
أسرع مما يمده به الهواء والتركيب الضوئي.
\textbf{23.} $2.80 \times 10^{-6}/(4 \times 96\,485) \times 3600 = \qty{2.6e-8}{mol}$
في الساعة: مهمل أمام أكسجين لتر واحد من الماء
(\qty{1.8e-4}{mol}).
\textbf{24.} تتغير مع درجة الحرارة ذوبانية \ce{O2} (\omterm{prop:b2:chemical-potential:henry}{ثابت هنري}) والانتشار
عبر الغشاء معًا: فلا تعود الحساسية المعيّنة عند
\qty{25}{\degreeCelsius} صالحة.
\textbf{25.} تحوي العيّنة $\boldsymbol{\approx \qty{5.9}{mg/L}}$ من الأكسجين
المذاب.
\end{solution}
